\documentclass[11pt]{article}
\usepackage[T1]{fontenc}
\usepackage{lmodern}
\usepackage{amsmath,amssymb,amsthm,mathtools}
\usepackage[left=1in,right=1in,top=0.85in,bottom=0.85in]{geometry}
\usepackage[colorlinks=true,linkcolor=blue,citecolor=blue,urlcolor=blue]{hyperref}
\usepackage{microtype}

\newtheorem{theorem}{Theorem}
\newtheorem{lemma}[theorem]{Lemma}
\newtheorem{remark}[theorem]{Remark}
\newcommand{\norm}[1]{\left\lVert #1\right\rVert}
\newcommand{\BX}{B_X}
\newcommand{\SX}{S_X}

\title{Local convexity of nonlinear self-maps characterizes quadratic convexity\\
\large A candidate full solution to Problem 6.1 of arXiv:1205.3412}
\author{Automated mathematics discovery run \texttt{fa\_banach\_001}\\
Lane 5, candidate for human review}
\date{21 July 2026}

\begin{document}
\maketitle

\begin{center}
\fbox{\parbox{0.92\textwidth}{\textbf{Status: candidate full solution, likely valid.}
The property in Problem 6.1 forces the given norm to have modulus of
convexity of power type $2$.  Consequently it also forces superreflexivity.
Together with the forward theorem in the source paper, this gives an
equivalence.}}
\end{center}

\section{The question and the answer}

For a Banach space $X$, Banakh, Banakh, Plichko and Prykarpatsky define
\[
 \operatorname{conv}_2(X)
 =\inf_{\substack{x,y\in \BX\\x\ne y}}
   \frac{1-\norm{(x+y)/2}}{\norm{x-y}^2}.
\]
This number is positive exactly when the given norm has modulus of
convexity of power type $2$.  Their main theorem says that if
$\operatorname{conv}_2(X)>0$, every Lipschitz-open, second-order Lipschitz
map from an open subset of $X$ to any Banach space is uniformly locally
convex \cite{BBPP}.

Problem 6.1 of the source asks for the converse when the codomain is $X$:
if every Lipschitz-open, second-order Lipschitz map $f:U\to X$ is locally
convex, must $X$ have modulus of convexity of power type $2$?  Must it be
(super)reflexive?  We answer both questions affirmatively.

All arguments below are over the real scalar field.  For a complex Banach
space we apply them to its underlying real Banach space.  This loses
nothing: the hypotheses in the problem are metric, and convexity is real
convexity.

\begin{theorem}\label{thm:main}
For a Banach space $X$, the following are equivalent.
\begin{enumerate}
\item $\operatorname{conv}_2(X)>0$.
\item Every Lipschitz-open, second-order Lipschitz map $f:U\to X$, with
      $U\subset X$ open, is locally convex.
\end{enumerate}
If these conditions hold, then $X$ is superreflexive (and hence reflexive).
\end{theorem}

The implication $(1)\Rightarrow(2)$ is the main theorem of the source
paper.  We prove the converse by contrapositive.  The counterexample will be
a quadratic perturbation of the identity on the unit ball.

\section{Proof intuition}

When $\operatorname{conv}_2(X)=0$, the unit ball has chords whose midpoint
lies only $o(\bigl(\text{chord length}\bigr)^2)$ below a supporting
hyperplane.  A quadratic term of the natural chord-squared size can therefore
push the preimage of an image midpoint outside a ball even when both endpoint
preimages remain inside.

We realize these pushes simultaneously with a norm-convergent bilinear
series.  The possible rotation of the supporting functionals is the one
nontrivial selection issue.  Either some fixed output direction is seen by
arbitrarily flat supports, or sufficiently flat supports become uniformly
small on every prescribed finite set of directions.  The first alternative
gives a rank-one range; the second gives an almost diagonal series whose
distinguished term dominates its past and tail.  Finally, making the
quadratic perturbation small compared with the identity guarantees
Lipschitz openness and controls its inverse displacement.

\section{Small quadratic perturbations are admissible}

\begin{lemma}\label{lem:quadratic}
Let $B:X\times X\to X$ be a continuous symmetric real-bilinear map with
$\norm{B}\leq\kappa$, where $0<2\kappa<1$.  On $U=B_1(0)$ put
\[
                     f(x)=x+B(x,x).
\]
Then $f$ is injective, Lipschitz-open, and second-order Lipschitz.  More
precisely, with $L=2\kappa$,
\begin{align}
 \norm{f(x)-f(y)}&\geq(1-L)\norm{x-y} \qquad(x,y\in U),\label{eq:lower}\\
 \operatorname{Lip}_o(f)&\geq1-L,\label{eq:open}\\
 \operatorname{Lip}_2(f)&\leq2\kappa.\label{eq:lip2}
\end{align}
\end{lemma}

\begin{proof}
Write $Q(x)=B(x,x)$.  For $x,y\in U$,
\[
 Q(x)-Q(y)=B(x-y,x)+B(y,x-y),
\]
so $\norm{Q(x)-Q(y)}\leq L\norm{x-y}$.  This proves
\eqref{eq:lower}, hence injectivity.

Suppose $B_r(x)\subset U$ and
$\norm{w-f(x)}<(1-L)r$.  Choose $s<r$ so that
$\norm{w-f(x)}<(1-L)s$.  On the complete closed ball
$\overline B_s(x)$, the map
\[
                         T(z)=w-Q(z)
\]
is an $L$-contraction and maps the ball strictly into itself, since
\[
 \norm{T(z)-x}
 \leq\norm{w-f(x)}+\norm{Q(x)-Q(z)}< (1-L)s+Ls=s.
\]
Its fixed point belongs to $B_s(x)$ and satisfies $f(z)=w$.  This proves
\eqref{eq:open}.  Finally,
\[
 f(x+h)-2f(x)+f(x-h)=2B(h,h),
\]
which gives \eqref{eq:lip2} with the definition used in the source.
\end{proof}

\section{A bilinear form detecting arbitrarily flat midpoints}

Assume henceforth that $\operatorname{conv}_2(X)=0$.  For
$x,y\in\BX$, set
\[
             z=\frac{x+y}{2},\qquad h=\frac{x-y}{2},\qquad
             d=1-\norm{z}.
\]
Call $(x,y,\sigma)$ a \emph{flat triple} if $x\ne y$, $z\ne0$,
$\sigma\in S_{X^*}$, and $\sigma(z)=\norm{z}$.  Its flatness ratio is
$d/\norm{h}^2$.  Because
\[
 \inf\frac{d}{\norm{h}^2}=4\operatorname{conv}_2(X)=0,
\]
flat triples with arbitrarily small ratio exist.

Set
\[
       \alpha_n=16^{-n},\qquad
       \kappa=\sum_{n=1}^\infty\alpha_n=\frac1{15},\qquad
       C=(1-2\kappa)^{-1}.
\]
We shall construct flat triples $(x_n,y_n,\sigma_n)$, vectors
$v_n\in\SX$, functionals $\varphi_n\in S_{X^*}$, positive numbers
$\beta_n$, and a symmetric bilinear map $B$ such that
\begin{equation}\label{eq:detection}
 \norm{B}\leq\kappa,
 \qquad
 \sigma_n\bigl(B(h_n,h_n)\bigr)
       \geq\beta_n\norm{h_n}^2,
\end{equation}
where $h_n=(x_n-y_n)/2$.  The only issue is that the norming functionals
$\sigma_n$ may rotate.  The following exhaustive dichotomy handles it.

\subsection*{Case I: one direction is seen infinitely often}

Suppose there are $v\in\SX$, $c>0$, and flat triples of arbitrarily small
ratio for which $|\sigma(v)|\geq c$.  By passing to one sign along a
sequence whose ratios tend to zero, and replacing $v$ by $-v$ if needed,
we may require $\sigma_n(v)\geq c$.  Put
\[
                         \beta_n=c\alpha_n.
\]
Choose $\varphi_n\in S_{X^*}$ with
$\varphi_n(h_n)=\norm{h_n}$, and define
\begin{equation}\label{eq:Bcase1}
 B(u,w)=\sum_{m=1}^\infty
       \alpha_m\varphi_m(u)\varphi_m(w)v.
\end{equation}
The series converges in operator norm and $\norm{B}\leq\kappa$.  Moreover,
all its scalar summands are nonnegative, so
\[
 \sigma_n(B(h_n,h_n))
 =\sigma_n(v)\sum_m\alpha_m\varphi_m(h_n)^2
 \geq c\alpha_n\norm{h_n}^2.
\]

\subsection*{Case II: every fixed direction eventually disappears}

Suppose Case I fails.  Then for every $v\in\SX$ and every $c>0$ there is a
threshold $\theta>0$ such that
\begin{equation}\label{eq:disappear}
 \frac{1-\norm{z}}{\norm{h}^2}<\theta
 \quad\Longrightarrow\quad |\sigma(v)|<c
\end{equation}
for every flat triple $(x,y,\sigma)$.  The same assertion holds
simultaneously for any finite list of directions, by taking the minimum of
the corresponding thresholds.

We choose the triples inductively.  Once $v_1,\ldots,v_{n-1}$ have been
chosen, use \eqref{eq:disappear} and then choose the next triple sufficiently
flat that
\begin{equation}\label{eq:pastsmall}
                 \sum_{m<n}\alpha_m|\sigma_n(v_m)|
                 \leq\frac{\alpha_n}{4}.
\end{equation}
Set
\[
 v_n=\frac{z_n}{\norm{z_n}},\qquad
 \beta_n=\frac{\alpha_n}{2},
\]
so $\sigma_n(v_n)=1$.  Again take
$\varphi_n\in S_{X^*}$ with
$\varphi_n(h_n)=\norm{h_n}$, and now define
\begin{equation}\label{eq:Bcase2}
 B(u,w)=\sum_{m=1}^\infty
       \alpha_m\varphi_m(u)\varphi_m(w)v_m.
\end{equation}
Then $\norm{B}\leq\kappa$.  Since
\[
                  \sum_{m>n}\alpha_m=\frac{\alpha_n}{15},
\]
the distinguished $n$th term, \eqref{eq:pastsmall}, and the tail estimate
give
\begin{align*}
 \sigma_n(B(h_n,h_n))
 &\geq \alpha_n\norm{h_n}^2
   -\norm{h_n}^2\sum_{m<n}\alpha_m|\sigma_n(v_m)|
   -\norm{h_n}^2\sum_{m>n}\alpha_m\\
 &\geq\left(1-\frac14-\frac1{15}\right)
       \alpha_n\norm{h_n}^2
 \geq\frac{\alpha_n}{2}\norm{h_n}^2.
\end{align*}
Thus \eqref{eq:detection} holds in both cases.

\section{The images of arbitrarily small balls are nonconvex}

There is a harmless order-of-choice point in the preceding construction
which we now make explicit.  In either case, before choosing the $n$th flat
triple choose $\varepsilon_n>0$ so that $\varepsilon_n\downarrow0$,
$\varepsilon_n<1/4$, and
\begin{equation}\label{eq:epssmall}
 2C\kappa^2\varepsilon_n+C^2\kappa^3\varepsilon_n^2
 \leq\frac{\beta_n}{2}.
\end{equation}
Then choose the triple sufficiently flat that, in addition to the relevant
case condition,
\begin{equation}\label{eq:flatneeded}
 d_n<\frac14\varepsilon_n\beta_n\norm{h_n}^2.
\end{equation}
This is possible in Case I by its definition, and in Case II by
\eqref{eq:disappear}.  Thus the $B$ just constructed satisfies all the
preceding estimates for these choices.

Apply Lemma~\ref{lem:quadratic} to this $B$ and let
\[
                     f(x)=x+B(x,x),\qquad x\in B_1(0).
\]
The map is Lipschitz-open and second-order Lipschitz.  Consider
\[
 a_n=\varepsilon_nx_n,\qquad c_n=\varepsilon_ny_n.
\]
The midpoint of their images is
\begin{align}
 w_n&=\frac{f(a_n)+f(c_n)}2\notag\\
    &=f(\varepsilon_nz_n)
      +\varepsilon_n^2B(h_n,h_n),\label{eq:image-midpoint}
\end{align}
where bilinearity was used in the second line.

Indeed $B_{1/2}(\varepsilon_nz_n)\subset B_1(0)$, and
\[
 \norm{w_n-f(\varepsilon_nz_n)}
 \leq\kappa\varepsilon_n^2
 <\frac{1-2\kappa}{2}.
\]
Thus Lipschitz openness gives an $r_n\in B_1(0)$ with $f(r_n)=w_n$;
injectivity makes it unique.  Write
\[
                         r_n=\varepsilon_nz_n+u_n.
\]
The lower Lipschitz estimate \eqref{eq:lower} and
\eqref{eq:image-midpoint} imply
\begin{equation}\label{eq:ubound}
 \norm{u_n}
 \leq C\varepsilon_n^2\norm{B(h_n,h_n)}
 \leq C\kappa\varepsilon_n^2\norm{h_n}^2.
\end{equation}
Expanding the equation $f(r_n)=w_n$ gives
\[
 u_n+2\varepsilon_nB(z_n,u_n)+B(u_n,u_n)
       =\varepsilon_n^2B(h_n,h_n).
\]
Consequently, using $\norm{z_n},\norm{h_n}\leq1$ and
\eqref{eq:ubound},
\begin{align*}
 \norm{u_n-\varepsilon_n^2B(h_n,h_n)}
 &\leq2\kappa\varepsilon_n\norm{u_n}
       +\kappa\norm{u_n}^2\\
 &\leq \varepsilon_n^2\norm{h_n}^2
       \bigl(2C\kappa^2\varepsilon_n
              +C^2\kappa^3\varepsilon_n^2\bigr)\\
 &\leq\frac12\varepsilon_n^2\beta_n\norm{h_n}^2.
\end{align*}
Together with \eqref{eq:detection}, this yields
\begin{equation}\label{eq:supportu}
              \sigma_n(u_n)
              \geq\frac12\varepsilon_n^2
                     \beta_n\norm{h_n}^2.
\end{equation}

Now put
\[
\rho_n=\varepsilon_n
       +\frac14\varepsilon_n^2\beta_n\norm{h_n}^2.
\]
Here $\beta_n\leq\alpha_1=1/16$, so $\rho_n<1$ and this ball lies in the
domain of $f$.  Both $a_n$ and $c_n$ lie in the open ball
$B_{\rho_n}(0)$, while
\eqref{eq:flatneeded} and \eqref{eq:supportu} give
\begin{align*}
 \norm{r_n}
 &\geq\sigma_n(r_n)\\
 &=\varepsilon_n(1-d_n)+\sigma_n(u_n)\\
 &>\varepsilon_n
   -\frac14\varepsilon_n^2\beta_n\norm{h_n}^2
   +\frac12\varepsilon_n^2\beta_n\norm{h_n}^2
 =\rho_n.
\end{align*}
Thus $r_n\notin B_{\rho_n}(0)$.  Since $f$ is injective,
$w_n$ has no other preimage in that ball.  But $w_n$ is the midpoint of
$f(a_n)$ and $f(c_n)$, so
\[
                         f(B_{\rho_n}(0))
\]
is not convex.  As $\rho_n\to0$, the map $f$ is not locally convex at the
origin.  This proves the contrapositive of $(2)\Rightarrow(1)$.

Finally, $\operatorname{conv}_2(X)>0$ makes the given norm uniformly convex.
Hence $X$ is superreflexive, and in particular reflexive.  This completes
the proof of Theorem~\ref{thm:main}.

\section{Scope and novelty check}

The construction answers both clauses of Problem 6.1 and does not require
differentiability beyond the source's second-order Lipschitz hypothesis.
It also strengthens the qualitative conclusion: failure of quadratic
convexity is witnessed by one explicit quadratic perturbation of the
identity whose images of arbitrarily small balls centered at zero are
nonconvex.

Cheap indexes for the present run and a bounded title/keyword search found
the source paper and later papers discussing local convexity, but no prior
proof of this converse.  This is not an exhaustive novelty certification.
The flat-triple dichotomy and the nonconvexity calculation should receive
independent specialist review.

\begin{thebibliography}{9}
\bibitem{BBPP}
I.~Banakh, T.~Banakh, A.~Plichko, and A.~Prykarpatsky,
\emph{On local convexity of nonlinear mappings between Banach spaces},
arXiv:1205.3412 (2012),
\url{https://arxiv.org/abs/1205.3412}.
\end{thebibliography}

\end{document}
